
\begin{exercise}[Lipschitz continuous bump functions \Coffeecup\Coffeecup]
    \label{ex: lipschitz bump functions}
    Let \((\domain, \metric)\) be a metric space and \(A \subseteq \domain\). 
    Show that
    \begin{enumerate}[label=(\alph*)]
        \item \(x\mapsto \metric(x,A)\coloneq \inf_{z \in A} \metric(x,z)\) is \(1\)-Lipschitz continuous.
        \begin{solution}
        By the triangle inequality we have for \(x,y\in \domain\)
        \[
            \metric(x,A)
            =\inf_{z\in A} \metric(x,z)
            \le \inf_{z\in A} \bigl(\metric(x,y) + \metric(y,z)\bigr)
            = \metric(x,y) + \metric(y,A).
        \]
        with reordering this implies
        \[
            \metric(x,A) - \metric(y,A) \le \metric(x,y)
        \]
        By symmetry we also have  \(\metric(y,A) - \metric(x,A) \le \metric(x,y)\)
        and thereby the claim.
        \end{solution}

        \item \(x\mapsto x_+ = \max\set{x,0}\) is \(1\)-Lipschitz continuous.
        \begin{solution}
            Choose \(x,y\in \real\), without loss of generality
            assume \(x \ge y\). If \(y\le x\le 0\), then \(x_+ = y_+ = 0\) and thus
            \(
                \abs{x_+ - y_+} = 0 \le \abs{x-y}. 
            \)

            If \(x\ge 0\), then
            \(
                \abs{x_+ - y_+} = x - y_+ \le x-y = \abs{x-y}.
            \)
        \end{solution}

        \item If \(f\) is \(L_1\)-Lipschitz continuous and \(g\) is \(L_2\)-Lipschitz continuous, then
        \(f\circ g\) is \(L_1L_2\)-Lipschitz continuous.
        \begin{solution}
            We have for any \(x,y\in \domain\)
            \[
                \abs{f(g(x)) - f(g(y))}
                \le L_1 \abs{g(x) - g(y)}
                \le L_1L_2 \abs{x-y}.
                \qedhere
            \]
        \end{solution}


        \item the bump function \(\ind{A}^\epsilon(x) = (1 - \frac{\metric(x,A)}\epsilon)_+\)
        is \(\frac1\epsilon\)-Lipschitz continuous.
        \begin{solution}
            \(f_1(x) = \metric(x,A)\) is \(1\)-Lipschitz continuous,
            \(f_2(x) = 1 - \frac{x}{\epsilon}\) is \(\frac1\epsilon\)-Lipschitz continuous
            as a Linear function, and \(f_3(x) = x_+\) is \(1\)-Lipschitz continuous. Consequently
            \[
                f_3 \circ f_2 \circ f_1(x) = \ind{A}^\epsilon(x)
            \]
            is \(\frac1\epsilon\)-Lipschitz continuous by the previous part.
        \end{solution}
    \end{enumerate}
\end{exercise}
